\subsection{The converse with compact zero exclusion}\label{sec:compact-rh}

Here is a constructive version of the reverse implication. Let
$E(s)=(s-1)\zeta(s)$ be the pole-removed function defined by the
fractional-part integral, including $E(1)=1$. Consider the arithmetic
condition
\[
 \text{for every } m,T\ge1\ \text{there is } b\ge1\ \text{such that }
       |E(s)|\ge 1/b
\]
on $1/2+1/m\le\operatorname{Re}s\le4$, $|\operatorname{Im}s|\le T$.
It suffices to quantify over rational complex points and approximation
precisions because the derivative bound from the integral supplies a PR
continuity modulus extending the estimate to all names in a rectangle.
For instance the lower bound means that every disk of radius
$2^{-j}$ about a returned approximant has center of magnitude
at least $1/b-2^{-j}$. We call this condition compact zero exclusion.
Beyond $\operatorname{Re}s=4$ the Euler product already gives uniform
separation. Classically, the usual functional-equation symmetry and
compactness identify zero exclusion on the right of the critical line
with RH. The present equivalence is between two statements of arithmetic.

\begin{theorem}\label{thm:compact-rh}
Compact zero exclusion is equivalent in the stated iPRA interpretation
to the family $\forall r\ge1\ \exists C\ge1\ \mathcal B(r,C)$.
In the reverse direction the bound at exponent $1/2+1/(2r)$ has an
explicit constant, independent of the compact separation witnesses $b$.
\end{theorem}
\begin{proof}
First assume the growth family. In the proof of
Theorem~\ref{thm:growth} choose $r\ge m$, so that
$\alpha_r\le1/2+1/(2m)$. Then
$|D(s)|\le2mC(T+4)$ on the indicated rectangle.
When $|s-1|\ge1/8$ use $D(s)E(s)=s-1$ to get
$|E(s)|\ge1/(16mC(T+4))$. When $|s-1|\le1/4$, the direct integral
estimate in Section~\ref{sec:logrates} gives $|E(s)|\ge1/3$.
The regions overlap, giving a single admissible $b$.

For the converse fix positive rationals $\gamma,\kappa\le1/4$ and
$T\ge4$. We first prove
\begin{equation}\label{eq:subpower}
 |D(s)|\le P_{\gamma,\kappa}(T+5)^\kappa,
 \qquad 1/2+\gamma\le\operatorname{Re}s\le2,\quad |\operatorname{Im}s|\le T,
 \qquad D(s)=\frac{s-1}{E(s)} .
\end{equation}
For a given height $t$ center a disk at $z_0=2+it$, of radius
$R=3/2-\gamma/2$. On it
$H(z)=E(z)/(z+1)$ is separated from zero using the witness at
$m=\lceil4/\gamma\rceil$ and height $\lceil T\rceil+3$. This covers
a neighborhood of the disk. The fractional-part integral and
$\operatorname{Re}z\ge1/2$ give $|H(z)|\le4(T+5)^2$:
use $|z|\le T+4$, $|z-1|\le T+5$, $|z+1|\ge1$, and
$\int_1^\infty u^{-\operatorname{Re}z-1}du\le2$.
Define a logarithm $L(z)$ of $H(z)$ by integrating $H'/H$ from $z_0$
on straight segments. Start it with the branch on the inner disk
$|z-z_0|\le1/2$ given by
\[
    \sum_{p,k\ge1}\frac{p^{-kz}}{k}
       -\sum_{k\ge1}\frac1k\left(\frac{2}{z+1}\right)^k,
       \qquad p\text{ prime}.
\]
The first sum has absolute majorant $\sum_{n\ge2}n^{-3/2}\le2$ by
distinctness of prime powers, and the second has majorant
$\sum_{k\ge1}(4/5)^k=4$. Their exponentials multiply to $H$, by the
Euler product and $1-2/(z+1)=(z-1)/(z+1)$.
The straight-segment primitive of an analytic $g$ is
$(z-z_0)\int_0^1g(z_0+v(z-z_0))\,dv$. Differentiation integrates
$\frac{d}{dv}[v g(z_0+v(z-z_0))]$, giving $g(z)$.
Thus $L'=H'/H$, and differentiating $H e^{-L}$ shows it is $1$.
On the inner disk $L$ agrees with the displayed branch by its
derivative and initial value. In particular $|L|\le6$ there, and on
the large disk $\operatorname{Re}L\le B=\log4+2\log(T+5)$.

Let $a_j$ be the Taylor coefficients of $L$ at $z_0$. On a circle of
radius $d\le R$ they are
\[
 a_jd^j=\frac1{2\pi}\int_0^{2\pi}
                L(z_0+d e^{i\theta})e^{-ij\theta}\,d\theta .
\]
For completeness, the circle integral of an analytic function vanishes
by the straight-segment primitive. Apply this to
$(L(z)-L(w))/(z-w)$, continued at $w$ by the segment integral of
$L'$ from $w$ to $z$ with parameter in $[0,1]$. For $|w-z_0|<d$
expand $1/(z-w)$ about $z_0$ on the circle by its geometric series.
Its normalized integral is $1$. The resulting Cauchy formula gives
the displayed coefficients and their convergent expansion at $w$.
In particular coefficients obtained from both circles agree
(alternatively differentiate the series at the center).
For $j\ge1$ the mean of $L(z_0+R e^{i\theta})e^{ij\theta}$ is zero,
by integrating $L(z)(z-z_0)^{j-1}$ around the circle.
Conjugating shows that the displayed mean for $a_jR^j$ equals
twice the mean of $\operatorname{Re}L\,e^{-ij\theta}$.
Subtract the constant $B$ in that last mean and use
$B-\operatorname{Re} L\ge0$, whose mean is $B-\operatorname{Re}a_0$.
Together with the inner circle this gives
\[
       |a_j|\le 6\,2^j\quad(j\ge0),\qquad
       |a_j|R^j\le2(B+6)\quad(j\ge1).
\]
Circle quadrature here uses rational parameter meshes, the named
sine and cosine, and their derivative bounds. A supplied $b$ controls
the denominators in the segment primitive. It suffices even when
it is very large.

At radius $r_0=3/2-\gamma$, write $d_0=1-r_0/R\ge\gamma/3$,
and choose $J=\lceil200/(\kappa\gamma^2)\rceil$. Through index
$J$ the sum of the absolute terms is at most $3^{J+3}$.
The remaining tail is at most
\[
   2(B+6)\frac{(1-d_0)^{J+1}}{d_0}
    \le \frac{162(1+\log(T+5))}{(J+1)\gamma^2}
    \le \kappa(1+\log(T+5)).
\]
Here $(1-d)^n\le1/(1+nd)$ for $0<d<1$ follows by the binomial
inequality. Since $|D|=|(z-1)/(z+1)|\,|e^{-L}|\le e^{|L|}$,
\eqref{eq:subpower} holds, for example with
$P_{\gamma,\kappa}=3^{3^{J+3}+1}$. Every point at the fixed height
in the stated strip is within $r_0$ of $z_0$. The separation witness
was needed for the existence and quadrature of the logarithm. The
bound on its Taylor tail used only its real-part upper bound.

Finally put $e=1/(2r)$, $\gamma=\kappa=e/4$, and take $T=N\ge4$.
At $x=N+1/2$, the half-integer Perron bound and the rectangle shift
of Section~\ref{sec:logrates} apply with left edge
$a=1/2+\gamma$. Use compact zero exclusion for the rectangle in
the proof that its boundary integral is zero. The truncation error
is at most $300xL/N$, where $L$ is the bit length of $2N+1$.
On each horizontal edge use $|s|\ge N$, $x^c\le2x$ as before.
On the left edge use $\int_{-N}^{N}|a+it|^{-1}dt\le4+2\log N$.
Since $x\le2N$, $(N+5)^\kappa\le6N^\kappa$ and
$L\le4(1+\log N)$, these three contributions together are bounded by
$4000 P_{\gamma,\kappa} (1+\log N)N^{a+\kappa}$.
Now $1+\log N\le(1+1/\gamma)N^\gamma$, and
$a+\kappa+\gamma\le1/2+e$. Thus
$C=\lceil4000 P_{\gamma,\kappa}(1+1/\gamma)\rceil$ works for
this $r$, including the three smaller positive integers by their
trivial bounds. Raising to the integer power $2r$ gives the
predicate $\mathcal B(r,C)$. In deriving the bound at each $N$ it
suffices to take the one separation witness for the corresponding
finite rectangle. There is no need for a choice function for these
witnesses.
\end{proof}
